\section{Things have always been pretty bad}

\begin{quotex}
the suffering inside the Church ``would have originated not from external enemies, but from within the same Church" …
\flright{\textsc{Pope Benedict XVI}}

\end{quotex}
The title was motivated by the idea that, as bad as things seem today, things were bad in yesteryear. So I used the
story of Abraham from Boccaccio's \emph{Decameron} to show the life of Rome in the 14th century. I
followed it up with a couple of personal reflections, to which someone took objection. Her point it that it was an
impious attack on the Church and harmful to the faithful. I omitted those reflections then, but have restored them at
the end.

The issue is really about the nature of authority and secrecy. Criticising certain individuals for public faults is not
an attack on the ideals of hierarchy and authority. After recent events, however, I can understand that point of view.
Those in control of an organisation don't want their faults to be exposed, since it would
undermine their credibility. This is more acute if the organisation provides an income. It certainly does not benefit
the consumers of the organisation's teachings. Or ``the rabble", as this person called.

Jesus warned us about wolves in sheep's clothing. It is not a fault to point out the wolves,
especially if the shepherds are unwilling to do so.

\paragraph{Hermetic Groups}
By their nature, hermetic groups form and reform as the need arises. Attempts to solidify them in an organisation with
grades, levels, etc. seldom last. When the do, the passage through the grades becomes perfunctory, and not due to any
special knowledge.

Hierarchies arise naturally, those who have gone before, will help those who come after. There is no need for official
grades, as a person can judge his own progress. Besides, you really can't fool the others in the
group for very long. In our groups, we are interested in the ``real" person, i.e., someone who is beginning to transcend
his quotidian waking consciousness. Our so-called normal lives are of little interest, except when they provide the raw
material for awakening to the true I.

Money, sex, and power are dangers to any organisation, so the temptation needs to be avoided in the first place.
Recently a high-ranking Buddhist, close to the Dalai Lama, was accused of sexual exploitation. New Buddhist initiates
are told that sex is restricted to vaginal intercourse between married couples. Does that mean one rule for the
``rabble" and another for those in power? And did the Dalai Lama know?

That is why I reject the notions of anonymity and impersonality. Maybe at one time they served a purpose, but not now.
Strange and contradictory claims often follow an anonymous writer after his death. The same applies to those who kept
their personal lives hidden under the guise that only the teachings mattered. Not so. Ideas come alive only when they
are in someone's consciousness; mere abstractions are dead ideas, not the sign of a higher
calling.

I've been called vulgar, and of dragging the teachings through ``rot and mud". That is the talk of
LARPers, who role play as high initiates by putting on a certain appearance. The point, however, is to show that the
path is difficult so that others will not be discouraged. Dante and Boccaccio, who actually were high initiates, are
far more earthy than I am. Nor were they shy about criticising church officials when it was justified.

\paragraph{Secret Knowledge and Authority}
There is the idea the Hermetism is about learning some secret teachings. They are considered to be revealed by someone
held up as an authority figure. Even if you get a hold of the teaching, you are informed that you
won't understand it without the ``key" that the authority has. Or else, they hold out the promise
of an even higher teaching.

Of course, that is only to protect the organisation. They create a cult of true believers who are high above the
``rabble" and know teachings that are intended for ``future" humanity.

It is just the opposite. We don't need yet another revelation, since we barely understand the ones
we have. Hermetism it not about some secret teachings, but rather about self-knowledge. The teachings are only fingers
pointing at the moon; the real work is in observing their effect in one's own being.

\paragraph{Knowledge, Being, and Understanding}
The distinction between knowledge of dianoia and gnosis of episteme seems hard to grasp. For example, I was just shown
an article about the influence of the Golden Dawn, on the premise that it tells us something useful about Hermetism.
Far from it.

Another way to put it is to make the distinction between knowledge, in the sense of book or textual knowledge, and
understanding which transcends it. This is often represented as the ``aha" experience when you finally ``get it".
I'm sure many of you can relate to that.

Zen abounds in such stories. The student hears the koans. He may ``know" it, since he can repeat it verbatim, but he
doesn't ``get it". That is a sudden illumination, and there is no mechanical or algorithmic way to
achieve it.

I like to use the analogy of a joke. You may have had the experience of telling a joke, but the hearer
didn't get it. The hearer ``knows" the joke; he can write it down, but he never understood it.
Someone else, on the other hand, might let out a spontaneous laugh. That is analogous to understanding.

The goal is to raise one's level of being through the Hermetic exercises. Only in that way can the
knowledge be understood. There are four levels of being.

\begin{itemize}
\item \textbf{Subconscious}: The twilight state, mental processes that are not consciously recognized. 
\item \textbf{Waking consciousness}: What the natural man considers his ordinary state. It is the subjective experience
of the I, although most of the time he is mostly in a dream like state just below consciousness. 
\item \textbf{Self-consciousness}: This is the beginning of the awareness of the real I. This can come about through
concentration, which gathers the disparate elements of one's psyche together. In this state, one
observes oneself objectively. 
\item \textbf{Absolute consciousness} 
\end{itemize}
\paragraph{What kind of Training}
Last week, I listened to part of an interview with \textbf{Timothy Dolan}, the archbishop of New York. In response to
yet another sex abuse crisis, the Archbishop said that they now give training to those who will work with children.

Hearing that, I wanted to jump through the radio and demand to know: who the hell needs special training to know not to
treat children as sex toys? A fortiori, anyone who needs such training in the first place is not suitable for the job.

\paragraph{Blaming the Victim}
Most of the cases involve teenage boys, around 15 years or so, and not children. I've sometimes
wondered why these priests never got beat up. I doubt that my circle of friends at that age would have succumbed to any
such proposition.

By the time my sons were 16, they knew how to handle a gun, they lifted weights with me, and fought in the PAL boxing
programs. If they had ever been involved in such situation, it would most likely have been a really bad experience for
the priest.

The answer is that the priests know how to choose their victims. In the confessional, they hear who is confessing to
same sex attraction. Most likely he would then offer extra counselling, but that is speculation.



\flrightit{Posted on 2018-09-28 by Cologero }
